% =====================================================================
%  SPESIFIKASI ARSITEKTUR
%  Bahan penyusunan arsitektur penelitian
%  Hilma Amalia Putri — Informatika, Fakultas Teknik, Universitas Siliwangi
% =====================================================================

\documentclass[12pt,a4paper,oneside]{report}

\ifdefined\XeTeXversion
  \usepackage{fontspec}
  \IfFileExists{fonts/times.ttf}{
    \setmainfont{times.ttf}[Path=fonts/, BoldFont=timesbd.ttf,
      ItalicFont=timesi.ttf, BoldItalicFont=timesbi.ttf]
  }{\usepackage{mathptmx}}
\else
  \usepackage[T1]{fontenc}\usepackage[utf8]{inputenc}\usepackage{mathptmx}
\fi

\usepackage[english,indonesian]{babel}
\usepackage[a4paper,top=4cm,left=4cm,bottom=3cm,right=3cm]{geometry}
\usepackage{setspace}
\doublespacing
\usepackage{indentfirst}
\setlength{\parindent}{1.27cm}

\usepackage{etoolbox}
\makeatletter
\def\ps@isita{%
  \let\@mkboth\@gobbletwo
  \def\@oddhead{\hfil\normalfont\normalsize\thepage}%
  \def\@evenhead{\normalfont\normalsize\thepage\hfil}%
  \def\@oddfoot{}\def\@evenfoot{}}
\makeatother
\pagestyle{isita}

\usepackage{titlesec}
\renewcommand{\thechapter}{\arabic{chapter}}
\titleformat{\chapter}[display]{\centering\bfseries\normalsize}
  {\thechapter}{0pt}{\MakeUppercase}
\titlespacing*{\chapter}{0pt}{0pt}{1.5\baselineskip}
\titleformat{\section}{\bfseries\normalsize}{\thesection}{0.6em}{}
\titleformat{\subsection}{\bfseries\normalsize}{\thesubsection}{0.6em}{}

\usepackage{graphicx}
\usepackage{float}
\usepackage{booktabs}
\usepackage{tabularx}
\usepackage{array}
\newcolumntype{L}[1]{>{\raggedright\arraybackslash}p{#1}}
\renewcommand{\tabularxcolumn}[1]{>{\raggedright\arraybackslash}m{#1}}
\usepackage{caption}
\captionsetup{font=small,labelfont=bf,labelsep=period,
  justification=centering,singlelinecheck=false}
\captionsetup[table]{position=top,skip=0.5\baselineskip}
\AtBeginEnvironment{tabularx}{\singlespacing\footnotesize\setlength{\tabcolsep}{3.5pt}}
\counterwithin{table}{chapter}
\renewcommand{\thetable}{\thechapter.\arabic{table}}

\usepackage{enumitem}
\newlist{poin}{enumerate}{2}
\setlist[poin]{label=\alph*.,leftmargin=*,align=left,itemsep=0pt,topsep=0.25\baselineskip}
\newlist{urut}{enumerate}{2}
\setlist[urut]{label=\arabic*.,leftmargin=*,align=left,itemsep=0pt,topsep=0.25\baselineskip}

\usepackage[round,authoryear]{natbib}
\bibliographystyle{plainnat}
\newcommand{\parencite}[1]{\citep{#1}}
\newcommand{\textcite}[1]{\citet{#1}}
\usepackage[hidelinks]{hyperref}

\begin{document}

\begin{titlepage}
  \centering
  \vspace*{0.5cm}
  {\bfseries\normalsize SPESIFIKASI ARSITEKTUR\par}
  \vspace{0.4cm}
  {\normalsize Bahan penyusunan arsitektur penelitian\par}
  \vspace{1.2cm}
  {\bfseries\normalsize VERIFIKASI BERLAPIS DAN KOMPARASI MODEL\\[0.2cm]
   PADA TAMPILAN TUJUH SEGMEN\\[0.2cm] TIMBANGAN DIGITAL\par}
  \vspace{2cm}
  {\normalsize Hilma Amalia Putri\\[0.2cm] [NPM]\par}
  \vspace{2cm}
  {\bfseries\normalsize JURUSAN INFORMATIKA\\[0.2cm]
   FAKULTAS TEKNIK UNIVERSITAS SILIWANGI\\[0.2cm] TASIKMALAYA\par}
  \vspace{1cm}
  {\normalsize 2026\par}
\end{titlepage}

\pagenumbering{roman}
\pagestyle{plain}
\tableofcontents
\clearpage

\pagenumbering{arabic}
\pagestyle{isita}

% =====================================================================
\chapter{Ruang Lingkup, Aktor, dan Konteks}

\section{Apa yang Dimodelkan}

Sistem yang diarsiteki membaca panel tujuh segmen pada timbangan digital dari
kamera ponsel, lalu mengembalikan tiga hal sekaligus: bacaan angka, keputusan
layak atau tidak, dan alasan keputusan itu. Objek bacaan adalah tampilan LED dan
LCD pasif dengan 1--8 digit dan 0--2 titik desimal, sebagaimana tampilan yang
dipakai pada timbangan gudang.

Arsitektur ini mencakup lapisan model dan lapisan keputusan. Antarmuka pengguna,
basis data produksi, dan sinkronisasi belum termasuk, dan disebut sebagai tahap
lanjutan pada bagian batas.

\section{Aktor}

\begin{table}[H]
\caption{Aktor dan kebutuhannya}
\label{tab:aktor}
\begin{tabularx}{\textwidth}{@{}L{2.6cm} L{5.0cm} X@{}}
\toprule
\textbf{Aktor} & \textbf{Kebutuhan} & \textbf{Hak akses} \\
\midrule
Operator timbang & Membaca panel, mencatat hasil, menimbang ulang bila ditolak & Menjalankan pembacaan, melihat alasan penolakan \\
Administrator produk & Menetapkan rentang berat sah dan toleransi per produk & Mengubah konfigurasi produk \\
Peneliti & Menjalankan eksperimen, mengukur metrik, menyusun laporan & Akses penuh ke data uji, registry, dan harness \\
\bottomrule
\end{tabularx}
\end{table}

\section{Konteks Operasional}

\begin{poin}
  \item Satu transaksi adalah satu penimbangan: satu produk, satu panel, tiga
        sampai lima bingkai. Target waktu satu pembacaan di bawah dua detik pada
        ponsel kelas menengah.
  \item Lokasi gudang sering tanpa jaringan. Jalur runtime berjalan luring
        sepenuhnya; hanya mode pembanding model bahasa-visi yang memerlukan
        koneksi.
  \item Timbangan di lapangan tidak selalu baru. Segmen mati, layar buram, dan
        pantulan lampu gudang harus dianggap keadaan normal, bukan pengecualian.
        Pantulan dilaporkan sebagai penyebab kegagalan utama pada pendekatan
        detektor-pengenal \parencite{low2024detecting}.
\end{poin}

% =====================================================================
\chapter{Kebutuhan}

\section{Kebutuhan Fungsional}

\begin{table}[H]
\caption{Kebutuhan fungsional}
\label{tab:fungsional}
\begin{tabularx}{\textwidth}{@{}L{1.1cm} L{5.3cm} X@{}}
\toprule
\textbf{Kode} & \textbf{Kebutuhan} & \textbf{Kriteria terima} \\
\midrule
F1 & Mengambil 3--5 bingkai beserta metadata kondisi & Bingkai tersimpan dengan nomor urut dan cap waktu \\
F2 & Menemukan area panel dan menghitung jumlah slot & Slot terdeteksi pada 50 citra uji manual sebelum dipakai \\
F3 & Menjalankan enam pengenal pada citra yang sama & Setiap pengenal menghasilkan satu string dan satu skor \\
F4 & Menyeragamkan keluaran semua pengenal & Pemisah desimal, nol depan, dan spasi mengikuti satu berkas aturan \\
F5 & Memeriksa bacaan lewat empat lapis verifikasi & Setiap bacaan menghasilkan keputusan dan nama lapis yang memutuskan \\
F6 & Menyimpan jejak audit per transaksi & Citra, string mentah, keputusan, dan alasan tersimpan bersama \\
F7 & Mengelola konfigurasi produk & Rentang berat dan toleransi dapat diubah tanpa mengubah kode \\
F8 & Mengelola bobot dan ambang & Versi bobot, ambang, dan tanggal berlaku tercatat \\
F9 & Menghitung metrik dan menyusun laporan & Satu perintah menghasilkan tabel metrik dan kurva \\
\bottomrule
\end{tabularx}
\end{table}

\section{Kebutuhan Nonfungsional}

\begin{table}[H]
\caption{Kebutuhan nonfungsional}
\label{tab:nonfungsional}
\begin{tabularx}{\textwidth}{@{}L{1.1cm} L{4.6cm} X@{}}
\toprule
\textbf{Kode} & \textbf{Aspek} & \textbf{Batasan} \\
\midrule
N1 & Luring & Seluruh jalur runtime tanpa koneksi, kecuali mode pembanding \\
N2 & Latensi & p95 di bawah dua detik per transaksi pada perangkat sasaran \\
N3 & Ketertelusuran & Setiap keputusan dapat dipertanggungjawabkan sampai lapis dan ambangnya \\
N4 & Determinisme & Versi pustaka, bobot, dan benih acak dikunci sebelum eksperimen \\
N5 & Lisensi data & Kumpulan data akademik seperti UFPR-AMR tidak diredistribusi \parencite{laroca2019convolutional} \\
N6 & Privasi & Citra tidak keluar perangkat kecuali pada mode pembanding yang disetujui \\
\bottomrule
\end{tabularx}
\end{table}

% =====================================================================
\chapter{Komponen}

\section{Daftar Komponen}

Dua belas komponen berikut membentuk arsitektur. Kolom terakhir menyebut dasar
pemilihannya, bukan sekadar nama pustaka.

\begin{table}[H]
\caption{Komponen inti: akuisisi sampai pengenal}
\label{tab:komponen}
\begin{tabularx}{\textwidth}{@{}L{0.9cm} L{2.9cm} L{3.4cm} L{2.6cm} X@{}}
\toprule
\textbf{Kode} & \textbf{Komponen} & \textbf{Masukan, keluaran} & \textbf{Teknologi} & \textbf{Dasar} \\
\midrule
C1 & Akuisisi & Bingkai kamera, metadata kondisi & Kamera ponsel & Pembacaan dari kamera ponsel lazim pada alat ukur \parencite{kulkarni2021cnn} \\
C2 & Praproses dan slot & Citra, area panel, jumlah slot, batas slot & OpenCV, profil proyeksi & Proyeksi dipakai untuk memisah digit \parencite{jiang2026method} \\
C3a & Pengenal PP-OCRv6 & Area panel, string, skor & PaddleOCR, tiga tingkat & 1,5--34,5 juta parameter \parencite{zhang2026ppocrv6} \\
C3b & Pengenal PP-OCRv5 & Area panel, string, skor & PaddleOCR & 5 juta parameter \parencite{zhang2026ppocrv5} \\
C3c & Pengenal CRNN-CTC & Area panel, string, skor & PyTorch, dilatih sendiri & Arsitektur acuan domain \parencite{shi2017crnn} \\
C3d & Pengenal DBNet + PARSeq & Area panel, string, skor & PaddleOCR, bobot PARSeq & Terbaik pada data tujuh segmen \parencite{low2024detecting, bautista2022parseq} \\
C3e & Pengenal ssocr & Area panel, string & Perangkat lunak klasik & Pembanding tanpa pelatihan \parencite{auerswal_ssocr} \\
C3f & Pengenal bahasa-visi & Citra terenkode, string & Layanan daring & Unggul transkripsi, lemah skala \parencite{ega2026leveraging, lin2026measurebench} \\
\bottomrule
\end{tabularx}
\end{table}

\begin{table}[H]
\caption{Komponen keputusan dan pendukung}
\label{tab:komponen2}
\begin{tabularx}{\textwidth}{@{}L{0.9cm} L{2.9cm} L{3.4cm} L{2.6cm} X@{}}
\toprule
\textbf{Kode} & \textbf{Komponen} & \textbf{Masukan, keluaran} & \textbf{Teknologi} & \textbf{Dasar} \\
\midrule
C4 & Normalisasi keluaran & String mentah, string baku & Aturan eksplisit & Tanpa aturan bersama, perbandingan tidak sah \\
C5 & Harness metrik & Keluaran semua model, label & Python & Semua angka dari satu skrip \\
C6 & Pemilih model & Tabel akurasi dan biaya, model dasar & Aturan keputusan & Akurasi dan biaya diputuskan bersama \\
C7 & Verifikasi & String model dasar, konteks, keputusan & Aturan dan kalibrasi & Empat lapis pada Bab V \\
C8 & Mesin keputusan & Hasil verifikasi, keputusan beralasan & Aturan & Keluaran harus keputusan \\
C9 & Pencatat audit & Keputusan, catatan transaksi & Berkas JSON & Tanpa jejak, kesalahan tak tertelusur \\
C10 & Evaluator & Keputusan, label, metrik keandalan & Python & Keandalan diukur dari keputusan \\
C11 & Registry & Bobot, ambang, versi & Berkas dan manifes & Versi dikunci untuk reproduksi \\
C12 & Konfigurasi produk & Rentang berat, toleransi & Berkas konfigurasi & Lapis fisik butuh pengetahuan produk \\
\bottomrule
\end{tabularx}
\end{table}\section{Pemisahan yang Wajib Dijaga}

\begin{poin}
  \item C3a--C3f tidak boleh berbagi praproses yang berbeda. Semua menerima
        citra hasil C2 yang sama, karena perbedaan masukan akan mengubah arti
        perbandingan.
  \item C5 tidak boleh dipanggil oleh runtime. Ia berjalan setelah eksperimen,
        memakai data uji yang sudah dikunci.
  \item C10 tidak boleh membaca apa pun dari C7 selain keputusan akhir. Membaca
        skor internal akan membuat evaluasi mengukur niat, bukan hasil.
\end{poin}

% =====================================================================
\chapter{Kontrak Data}

\section{Keluaran Pengenal}

Setiap pengenal mengembalikan satu catatan dengan bentuk yang sama, agar C5
memperlakukan semuanya setara.

\begin{table}[H]
\caption{Kontrak keluaran setiap pengenal}
\label{tab:kontrak-pengenal}
\begin{tabularx}{\textwidth}{@{}L{2.4cm} L{1.8cm} L{3.2cm} X@{}}
\toprule
\textbf{Medan} & \textbf{Tipe} & \textbf{Contoh} & \textbf{Aturan} \\
\midrule
model & teks & ppocrv6\_small & Nama baku, tanpa spasi \\
teks & teks & 1006.5 & Titik sebagai pemisah desimal, tanpa spasi \\
skor & desimal & 0.9100 & Rata-rata peluang karakter, 4 angka di belakang koma \\
latensi\_ms & bilangan & 82 & Diukur pada perangkat sasaran, bukan mesin latih \\
parameter & bilangan & 7700000 & Jumlah parameter, untuk tabel biaya \\
versi & teks & 6.1.0 & Versi pustaka dan bobot \\
\bottomrule
\end{tabularx}
\end{table}

\section{Kontrak Keputusan}

\begin{table}[H]
\caption{Kontrak keputusan verifikasi}
\label{tab:kontrak-keputusan}
\begin{tabularx}{\textwidth}{@{}L{2.4cm} L{1.8cm} L{3.0cm} X@{}}
\toprule
\textbf{Medan} & \textbf{Tipe} & \textbf{Contoh} & \textbf{Aturan} \\
\midrule
bacaan & teks & 1006.5 & String setelah normalisasi \\
keputusan & pilihan & terima & Hanya terima, tolak, atau konfirmasi \\
lapis & teks & L1 & Lapis pertama yang memutuskan \\
alasan & teks & slot 6 dari 6 & Kalimat pendek yang dapat dibaca operator \\
skor\_terkalibrasi & desimal & 0.9300 & Hasil kalibrasi, bukan skor mentah \\
ambang & desimal & 0.9000 & Nilai ambang yang berlaku saat itu \\
waktu & teks & 2026-11-02T09:14:03 & Cap waktu ISO 8601 \\
\bottomrule
\end{tabularx}
\end{table}

\section{Berkas Label Data Uji}

\begin{table}[H]
\caption{Kolom berkas label}
\label{tab:label}
\begin{tabularx}{\textwidth}{@{}L{2.6cm} L{4.6cm} X@{}}
\toprule
\textbf{Kolom} & \textbf{Isi} & \textbf{Aturan} \\
\midrule
berkas & nama citra & Tanpa folder, unik \\
label & bacaan sebenarnya & Format sama dengan kontrak, titik desimal \\
jenis\_tampilan & led\_merah, led\_hijau, lcd\_pasif & Tiga nilai, tanpa variasi lain \\
kondisi & bersih, pantulan, miring, gelap & Satu nilai per citra \\
split & latih, validasi, kalibrasi, uji & Terkunci sebelum eksperimen \\
\bottomrule
\end{tabularx}
\end{table}

\section{Aturan Normalisasi}

\begin{poin}
  \item Koma dan titik sama-sama dibaca sebagai pemisah desimal, lalu ditulis
        ulang sebagai titik.
  \item Nol di depan dipertahankan bila ada di panel, karena nol itu bagian dari
        bacaan, bukan hiasan.
  \item Spasi, tanda kutip, dan satuan seperti "kg" dibuang sebelum dinilai.
  \item String kosong dan string lebih dari delapan karakter dihitung sebagai
        kegagalan, bukan dibuang dari penyebut.
\end{poin}

% =====================================================================
\chapter{Lapisan Verifikasi}

\section{Definisi Operasional}

\begin{table}[H]
\caption{Empat lapis verifikasi, aturannya, dan asal ambangnya}
\label{tab:lapis}
\begin{tabularx}{\textwidth}{@{}L{1.0cm} L{3.0cm} L{4.4cm} L{2.4cm} X@{}}
\toprule
\textbf{Lapis} & \textbf{Pertanyaan} & \textbf{Aturan operasional} & \textbf{Ambang} & \textbf{Asal} \\
\midrule
L1 & Apakah bentuknya masuk akal & Jumlah karakter tanpa titik sama dengan jumlah slot; titik desimal paling banyak satu dan berada di kandidat slot & N slot dari C2 & Profil proyeksi \parencite{jiang2026method} \\
L2 & Apakah nilainya masuk akal & Berat berada di antara batas produk; selisih terhadap median riwayat tidak melebihi toleransi & Batas produk dari C12 & Konfigurasi operator \\
L3 & Apakah bacaannya stabil & Tiga sampai lima bingkai; minimal empat dari lima identik & 4 dari 5 & Penyaringan antar bingkai \parencite{yenneti2023lightweight} \\
L4 & Seberapa yakin model & Skor setelah kalibrasi tidak kurang dari ambang & Dipilih dari data kalibrasi & Peluang CTC \parencite{graves2006ctc} \\
\bottomrule
\end{tabularx}
\end{table}

\section{Cara Menetapkan Ambang}

Ambang tidak boleh ditentukan dari data uji. Prosedur yang dipakai:

\begin{urut}
  \item Memisahkan data kalibrasi dari data uji sejak awal, dengan pembagian
        yang sama seperti data latih dan validasi.
  \item Mengkalibrasi skor pada data kalibrasi, lalu memeriksa kesesuaian antara
        skor dan frekuensi benar.
  \item Memilih ambang pada kurva \textit{risk--coverage} dari data kalibrasi,
        sesuai \textit{coverage} yang ditargetkan, misalnya 80 persen bacaan
        diterima otomatis.
  \item Mengunci ambang, mencatatnya di registry, lalu menjalankan data uji satu
        kali.
\end{urut}

\section{Urutan Pemeriksaan}

Pemeriksaan berjalan berurutan, dan lapis pertama yang menolak menghentikan
pemeriksaan berikutnya. Urutan ini disengaja: L1 paling murah dan paling sering
menangkap kesalahan bentuk, sedangkan L4 paling mahal karena memerlukan
kalibrasi dan tidak dapat menjelaskan alasan selain angka.

% =====================================================================
\chapter{Alur Keputusan}

\section{Keadaan}

\begin{table}[H]
\caption{Keadaan keputusan, pemicu, dan tindakan}
\label{tab:keadaan}
\begin{tabularx}{\textwidth}{@{}L{2.6cm} L{4.6cm} L{3.4cm} X@{}}
\toprule
\textbf{Keadaan} & \textbf{Pemicu} & \textbf{Tindakan sistem} & \textbf{Yang dilakukan operator} \\
\midrule
Diterima & Lolos L1--L4 & Catat bacaan, simpan jejak & Melanjutkan ke karung berikutnya \\
Ditolak & Gagal L1 atau L2 & Tampilkan alasan, jangan catat & Menimbang ulang atau memasukkan manual \\
Konfirmasi & Lolos L1 dan L2, gagal L3 atau L4 & Minta satu pengambilan tambahan & Mengarahkan kamera, memotret ulang \\
Batal & Tidak ada bingkai terbaca & Kembali ke keadaan awal & Memeriksa posisi kamera \\
\bottomrule
\end{tabularx}
\end{table}

\section{Penanganan Kegagalan Tiap Lapis}

\begin{poin}
  \item L1 gagal karena slot tidak terdeteksi: sistem memakai jumlah slot dari
        konfigurasi produk, dan menandai bahwa slot berasal dari konfigurasi,
        bukan dari citra.
  \item L2 gagal karena produk belum dipilih: sistem menolak menilai dan meminta
        operator memilih produk lebih dahulu.
  \item L3 gagal karena bacaan berubah-ubah: sistem menyimpan semua bingkai
        sebagai bukti, bukan hanya yang terakhir.
  \item L4 gagal karena skor rendah: sistem tetap menyimpan string mentah, agar
        kesalahan dapat dianalisis setelah kejadian.
\end{poin}

% =====================================================================
\chapter{Metrik dan Cara Menghitung}

\begin{table}[H]
\caption{Metrik, definisi, dan tempat pemakaiannya}
\label{tab:metrik}
\begin{tabularx}{\textwidth}{@{}L{2.6cm} L{6.2cm} L{3.0cm} X@{}}
\toprule
\textbf{Metrik} & \textbf{Definisi} & \textbf{Dipakai di} & \textbf{Satuan} \\
\midrule
FSR & Jumlah bacaan yang benar seluruhnya dibagi jumlah bacaan & Tahap model & persen \\
CER & Jumlah karakter salah dibagi jumlah karakter label & Tahap model & persen \\
Akurasi desimal & Bacaan benar pada posisi titik desimal & Tahap model & persen \\
Latensi p50 dan p95 & Waktu satu bacaan, diukur pada perangkat sasaran & Tahap biaya & milidetik \\
Kesalahan senyap & Bacaan salah yang diterima otomatis dibagi seluruh bacaan & Verifikasi & persen \\
Coverage & Bacaan yang diterima otomatis dibagi seluruh bacaan & Verifikasi & persen \\
Selective accuracy & Bacaan benar di antara yang diterima, dibagi yang diterima & Verifikasi & persen \\
AURC & Luas daerah di bawah kurva risiko terhadap coverage & Verifikasi & tanpa satuan \\
\bottomrule
\end{tabularx}
\end{table}

Angka tunggal tanpa rentang tidak dilaporkan. Setiap persentase disertai jumlah
sampel dan interval kepercayaan 95 persen dari 1.000 kali \textit{bootstrap}.

% =====================================================================
\chapter{Rancangan Eksperimen}

\begin{table}[H]
\caption{Eksperimen tahap model}
\label{tab:eks-a}
\begin{tabularx}{\textwidth}{@{}L{1.0cm} L{5.6cm} X@{}}
\toprule
\textbf{Kode} & \textbf{Konfigurasi} & \textbf{Pertanyaan yang dijawab} \\
\midrule
A1 & PP-OCRv6 tiny, small, medium tanpa adaptasi & Tingkat mana yang paling efisien pada perangkat sasaran? \\
A2 & PP-OCRv5 tanpa adaptasi & Apakah kenaikan generasi memberi manfaat pada domain ini? \\
A3 & CRNN-CTC dilatih pada data latih yang sama & Seberapa jauh model khusus tertinggal atau unggul? \\
A4 & DBNet + PARSeq & Apakah pipeline dua tahap lebih baik daripada satu model terpadu? \\
A5 & ssocr tanpa pelatihan & Apakah pembelajaran mendalam memang diperlukan? \\
A6 & Model bahasa-visi dengan prompt tetap & Berapa biaya untuk akurasi setara? \\
\bottomrule
\end{tabularx}
\end{table}

\begin{table}[H]
\caption{Eksperimen tahap verifikasi}
\label{tab:eks-b}
\begin{tabularx}{\textwidth}{@{}L{1.0cm} L{5.6cm} X@{}}
\toprule
\textbf{Kode} & \textbf{Konfigurasi} & \textbf{Pertanyaan yang dijawab} \\
\midrule
B1 & Model dasar tanpa verifikasi & Berapa kesalahan senyap pada keadaan dasar? \\
B2 & Ditambah L1 struktural & Berapa kesalahan bentuk tertangkap? \\
B3 & Ditambah L2 fisik & Tambahan manfaat dari konteks produk? \\
B4 & Ditambah L3 temporal & Tambahan manfaat dari antar bingkai? \\
B5 & Ditambah L4 kalibrasi & Di mana titik seimbang coverage dan akurasi? \\
\bottomrule
\end{tabularx}
\end{table}

\section{Uji Statistik}

\begin{poin}
  \item Dua model pada data uji yang sama diuji dengan McNemar, karena yang
        dibandingkan adalah benar atau salah pada citra yang sama.
  \item Selisih persentase dilaporkan dengan interval kepercayaan dari
        \textit{bootstrap}, bukan dari asumsi sebaran normal.
  \item Kesesuaian sebaran citra sintetis dan nyata diuji dengan Kolmogorov--
        Smirnov pada enam fitur: lebar digit relatif, rasio kontras, ketajaman,
        kemiringan, tingkat derau, dan jumlah slot.
\end{poin}

% =====================================================================
\chapter{Lingkungan Teknis}

\begin{table}[H]
\caption{Perangkat lunak dan perangkat keras}
\label{tab:lingkungan}
\begin{tabularx}{\textwidth}{@{}L{2.8cm} L{5.2cm} X@{}}
\toprule
\textbf{Bagian} & \textbf{Spesifikasi} & \textbf{Catatan} \\
\midrule
Bahasa & Python 3.13 & Satu versi untuk seluruh eksperimen \\
Pustaka pengenal & PaddleOCR untuk PP-OCRv5 dan PP-OCRv6 & Versi dikunci, bobot disimpan di registry \\
Pustaka lain & PyTorch, OpenCV, Pillow, NumPy & Versi dicatat pada manifes \\
Latin runtime & ONNX Runtime & Model diekspor ke ONNX agar tidak bergantung Paddle \\
Pelatihan & CPU kampus, atau GPU sewa bila tersedia & CRNN kecil, latihan dapat selesai dalam jam \\
Perangkat uji & Ponsel kelas menengah, tanpa koneksi & Tempat latensi diukur; bukan mesin latih \\
Penyimpanan & Berkas JSON per transaksi & Satu berkas per transaksi, nama mengandung cap waktu \\
\bottomrule
\end{tabularx}
\end{table}

% =====================================================================
\chapter{Risiko dan Mitigasi}

\begin{table}[H]
\caption{Risiko teknis dan mitigasinya}
\label{tab:risiko}
\begin{tabularx}{\textwidth}{@{}L{3.4cm} L{3.6cm} L{4.6cm} X@{}}
\toprule
\textbf{Risiko} & \textbf{Dampak} & \textbf{Mitigasi} & \textbf{Tanda peringatan} \\
\midrule
PP-OCRv6 sudah tinggi tanpa adaptasi & Kontribusi model melemah & Kontribusi digeser ke verifikasi dan biaya & FSR di atas 95 persen pada A1 \\
Model klasik mengungguli model besar & Klaim model terbaru goyah & Laporkan per kondisi, bukan rata-rata & Selisih besar pada kondisi bersih \\
Layanan bahasa-visi berubah & Hasil tidak dapat diulang & Kunci versi, simpan keluaran mentah & Format jawaban berubah saat diuji ulang \\
Gap sintetis dan nyata & Latihan tidak berbuah & Uji sebaran enam fitur sebelum latih & Uji KS menolak kesamaan \\
Slot gagal terdeteksi & L1 salah menolak & Fallback ke jumlah slot konfigurasi & Tingkat penolakan L1 di atas 30 persen \\
Label uji keliru & Semua metrik salah & Anotasi ganda 10 persen citra & Kesepakatan antar anotator di bawah 95 persen \\
\bottomrule
\end{tabularx}
\end{table}

% =====================================================================
\chapter{Peta Klaim ke Sumber}

Setiap angka pada dokumen ini berasal dari sumber berikut, bukan dari perkiraan.

\begin{table}[H]
\caption{Klaim yang dipakai dan sumbernya}
\label{tab:peta}
\begin{tabularx}{\textwidth}{@{}L{1.0cm} L{6.2cm} L{2.0cm} X@{}}
\toprule
\textbf{ID} & \textbf{Klaim} & \textbf{Angka} & \textbf{Keterangan} \\
\midrule
S01 & CRNN-CTC dengan pra-pemrosesan YOLO pada data logger tujuh segmen & 90,37 persen & 197 dari 218 citra, CER 3,6 persen \\
S02 & Model bahasa-visi pada perangkat ukur digital & 93,85 persen & Transkripsi digit, bukan pembacaan skala \\
S03 & Kumpulan data meter angka & 2.000 citra & Pembagian 800/400/800, 10.000 anotasi digit \\
S04 & Detektor dan pengenal terbaik pada data tujuh segmen & -- & DBNet dan PARSeq, gagal saat pantulan \\
S07 & Pembangkit citra sintetis tujuh segmen & 96,8 persen & Model digit dan desimal terpisah \\
S08 & Kepala inferensi digit pada video & 97,8 persen & Per bingkai, 30 bingkai per detik \\
S09 & Pendekatan klasik pada timbangan elektronik & 99,7 persen & 1.146 karakter, 20 milidetik \\
S13 & Tolok ukur alat ukur untuk model bahasa-visi & 30,2 persen & 2.442 citra, 26 jenis alat \\
S14 & PP-OCRv6 & 1,5 / 7,7 / 34,5 juta & 83,2 persen pada 15 kategori \\
S17 & YOLO dan CRNN pada layar tujuh segmen & 98 persen & Sekitar 800 foto lapangan \\
S22 & PP-OCRv5 & 5 juta parameter & Menyaingi model miliaran parameter \\
\bottomrule
\end{tabularx}
\end{table}

% =====================================================================
\chapter{Bahan Gambar Arsitektur}

Bagian ini menyebut gambar apa saja yang perlu dibuat, dan apa yang wajib ada di
dalamnya. Urutannya mengikuti tingkat rinci: dari konteks sampai keadaan.

\begin{table}[H]
\caption{Daftar gambar yang perlu dibuat}
\label{tab:gambar}
\begin{tabularx}{\textwidth}{@{}L{1.0cm} L{3.4cm} L{5.6cm} X@{}}
\toprule
\textbf{Kode} & \textbf{Gambar} & \textbf{Elemen wajib} & \textbf{Relasi wajib} \\
\midrule
D1 & Konteks & Operator, kamera, timbangan, sistem, registry & Siapa memberi masukan, siapa menerima keluaran \\
D2 & Komponen & C1 sampai C12 dengan kontrak singkat & Setiap panah diberi nama data yang berpindah \\
D3 & Alur runtime & Bingkai, praproses, enam model, verifikasi, keputusan & Urutan satu transaksi, termasuk cabang gagal \\
D4 & Sekuens & Satu transaksi berhasil dan satu ditolak & Waktu vertikal, pemanggilan antar komponen \\
D5 & Keadaan & Empat keadaan keputusan & Pemicu perpindahan dan tindakan operator \\
D6 & Penyiapan & Pembangkit, latihan, kalibrasi, ekspor, registry & Kapan tiap artefak dihasilkan \\
D7 & Evaluasi & Data uji terkunci, harness, metrik, laporan & Bahwa label tidak pernah masuk ke runtime \\
\bottomrule
\end{tabularx}
\end{table}

\section{Aturan Menggambar}

\begin{poin}
  \item Setiap panah diberi nama data, bukan kata kerja. "string dan skor" lebih
        baik daripada "memproses".
  \item Kotak komponen memuat nama dan satu keterangan dalam tanda kurung,
        misalnya "Enam model (kandidat 3)".
  \item Penyimpanan digambar sebagai silinder: registry, catatan transaksi,
        arsip hasil.
  \item Jalur yang berjalan di luar perangkat dipisahkan dari jalur di perangkat,
        karena latensi hanya bermakna pada perangkat sasaran.
  \item Warna hanya dipakai untuk membedakan tiga hal: data, proses, dan
        penyimpanan. Lebih dari itu mengalihkan perhatian dari isi.
\end{poin}

\clearpage
\bibliography{references}

\end{document}
